\chapter{The hyperbolic arc-length reconstruction engine}
\label{chap:hyperbolic_arclength}

% Chapter for Gluck/Hyperbolic/ArcLength.lean (aggregator) and its modules
% Gluck/Hyperbolic/ArcLength/{Ode, Closing, ConverseCap}.lean. All
% declarations are landed and sorry-free. The engine realizes
% genuinely-negative-curvature H^2 four-vertex profiles in the Poincare
% disk, unreachable by the convex-only flow of chap:spaceform_core. An
% earlier concrete gate-profile showcase (Gate/ForkA/ForkARobust and its
% hypothesis-free realization) was superseded by the general fork-A family
% chain of chap:hyperbolic_family and removed.

\section{Overview}

The tangent-angle flow engine of \cref{chap:spaceform_core} is
\emph{convex-only} for $\varepsilon=-1$: its trajectories have monotone tangent
angle and force the admissibility bracket
$\kappa-\varepsilon\langle z,ie^{i\theta}\rangle>0$, so negative geodesic
curvature is unreachable. This chapter builds the complementary engine:
parametrize by \emph{Euclidean arc length} $\sigma$ and drive the tangent angle
$\varphi$ by a first-order ODE,
\[
  z'(\sigma) = e^{i\varphi(\sigma)}, \qquad
  \varphi'(\sigma) =
  \frac{2\bigl(\kappa(\sigma) + \langle z(\sigma),
    i e^{i\varphi(\sigma)}\rangle\bigr)}{1 - \|z(\sigma)\|^{2}},
\]
a coupled system in $W=(z,\varphi)\in\{\|z\|<1\}\times\mathbb R$. The
denominator $1-\|z\|^2>0$ is the \emph{metric factor} of the Poincar\'e disk —
a confinement condition, \emph{not} an admissibility denominator — so the
system is well-posed for \emph{any} continuous $\kappa$, including
genuinely-negative values, and $\varphi$ may be non-monotone. Confinement
$\|z\|<1$ is not automatic (unit Euclidean speed can reach the boundary) and is
the crux estimate, handled relative to a bounded reference reconstruction by
$L^1$-Grönwall continuous dependence (\cref{lem:arc_trajectory_diff_bound},
\cref{lem:arc_trajectory_gronwall}).

Second asymmetry against the Euclidean engine of \cref{chap:arclength}:
H\textsuperscript{2} has \emph{no metric rescaling}, so the window length $L$
is \emph{co-constructed} with the profile. The converse closes on $[0,L]$ and
transports to the $2\pi$-convention only through the linear window
reparametrization $\chi(t)=(L/2\pi)t$; its conclusion is honestly
\emph{up-to-reparametrization} (\cref{thm:arclength_h2_converse},
\cref{lem:realizes_h2_of_reparam}). The explicit-window form
\cref{lem:arclength_h2_converse_at} exposes the window shift law
$\chi(t+2\pi)=\chi(t)+L$ — the conjugation datum the degree-one analysis of
\cref{chap:hyperbolic_converse} consumes to remove the reparametrization for
the concrete construction.

Three stages. \emph{ODE} (\texttt{Ode.lean}): truncated field,
Picard--Lindelöf flow, the $L^1$-Grönwall trajectory bounds, the
$\mathrm{Realizes}(-1)$ lemma, and the closed-form constant-curvature arc
model. \emph{Closing} (\texttt{Closing.lean}): the project-proved
Poincaré--Miranda theorem — the topological engine of the 2-D closing
shooting problems — and the two-arc quarter-period model endpoints with
their scalar closed-form reductions. \emph{Capstone}
(\texttt{ConverseCap.lean}): floor-glued periodic extension of a closing
window solution, chord-condition simplicity (reusing the Euclidean-in-disk
chord machinery of \cref{chap:simplicity} — embeddedness is a
$\mathbb C$-property, blind to the metric), and the arc-length converse.
The consumers — bicircle family, transport, turning, closing, simplicity and
the capstone realization — are the fork-A family chain of
\cref{chap:hyperbolic_family}.

\section{Declarations}

\begin{definition}\leanok
[H\texorpdfstring{$^2$}{2} arc-length angle speed]
  \label{def:arc_angle_speed}
  \lean{Gluck.Hyperbolic.arcAngleSpeed}
  \uses{def:spaceform_realizes}
  $\texttt{arcAngleSpeed}\;\kappa\,\sigma\,z\,\varphi =
  2\bigl(\kappa(\sigma)+\langle z, i e^{i\varphi}\rangle\bigr)/(1-\|z\|^{2})$:
  the algebraic solution of the $\mathrm{Realizes}(-1)$ gauge-speed relation
  $\tfrac{1-\|z\|^2}{2}\varphi'=\kappa+\langle z,ie^{i\varphi}\rangle$ at unit
  Euclidean speed. The denominator is the metric factor, positive on the disk
  for \emph{any} $\kappa$ — unlike the admissibility bracket of
  \cref{chap:spaceform_core}, no sign condition on $\kappa$ is imposed.
  Junk-value total function.
\end{definition}

\begin{definition}\leanok
[the truncated field and its flow]
  \label{def:arc_flow}
  \lean{Gluck.Hyperbolic.arcFlow}
  \uses{def:arc_angle_speed, def:spaceform_flow}
  The truncated reconstruction field (\texttt{arcField}) is
  $G_{\kappa,R}(\sigma,(z,\varphi)) = \bigl(e^{i\varphi},\;
  \texttt{arcAngleSpeed}\;\kappa\,\sigma\,(\mathrm{clamp}_R z)\,\varphi\bigr)$,
  with $\mathrm{clamp}_R$ (\texttt{clampBall}) the $1$-Lipschitz radial
  projection onto $\overline B(0,R)$: for $0\le R<1$, $|\kappa|\le M$ it is
  globally bounded and Lipschitz on $\mathbb C\times\mathbb R$
  (\texttt{arcField\_lipschitz}), and on $\|z\|\le R$ the clamp is inactive so
  it agrees with the true system (\texttt{truncatedArcAngleSpeed\_eq}).
  \texttt{arcFlow} is its chosen Picard--Lindelöf flow
  $\Psi_{\kappa,R,L,M,r_0}$: for $\|W_0\|\le r_0$, $\Psi(W_0,0)=W_0$ and
  $\Psi(W_0,\cdot)$ solves $W'=G_{\kappa,R}(\sigma,W)$ on $[0,L]$
  (\texttt{arcFlow\_spec}), uniquely (\texttt{arcFlow\_unique}, from the global
  space-Lipschitz constant). Coupled $(z,\varphi)$ analogue of the space-form
  flow of \cref{chap:spaceform_core}.
\end{definition}

\begin{lemma}\leanok
[$L^1$-Grönwall trajectory comparison]
  \label{lem:arc_trajectory_diff_bound}
  \lean{Gluck.Hyperbolic.arcTrajectory_diff_bound}
  \uses{def:arc_flow}
  For solutions $W$, $W_s$ of the $\kappa$- and $\kappa'$-arc-length ODEs on
  $[0,L]$, with $\mathrm{Lip}$ a Lipschitz witness for the $\kappa$-field,
  \[
    \|W(\sigma)-W_s(\sigma)\| \le \|W(0)-W_s(0)\|
    + \int_0^{\sigma}\Bigl(\mathrm{Lip}\,\|W-W_s\|
      + \tfrac{2}{1-R^2}\,|\kappa-\kappa'|\Bigr).
  \]
  The curvature enters only through the $L^1$ distance — the workhorse of the
  family chain's confinement and robustness estimates
  (\cref{chap:hyperbolic_family}).
\end{lemma}

\begin{proof}
\leanok
  FTC on $W-W_s$ plus the pointwise field bound
  $\|G_{\kappa,R}(\sigma,W)-G_{\kappa',R}(\sigma,W')\|\le
  \mathrm{Lip}\|W-W'\|+\tfrac{2}{1-R^2}|\kappa(\sigma)-\kappa'(\sigma)|$
  (\texttt{arcField\_sub\_le}).
\end{proof}

\begin{lemma}\leanok
[$L^1$-Grönwall trajectory bound, exponential form]
  \label{lem:arc_trajectory_gronwall}
  \lean{Gluck.Hyperbolic.arcTrajectory_gronwall}
  \uses{lem:arc_trajectory_diff_bound, def:arc_flow, lem:gronwall_L1_drive}
  For solutions $W$, $W_s$ of the $\kappa$- and $\kappa'$-arc-length ODEs on
  $[0,L]$,
  \[
    \|W(t)-W_s(t)\| \le e^{\mathrm{Lip}\cdot L}\Bigl(\|W(0)-W_s(0)\|
      + \tfrac{2}{1-R^2}\int_0^{L}|\kappa-\kappa'|\Bigr).
  \]
  Two trajectories with $L^1$-close curvatures and close starts stay close —
  the transport input of the fork-A family layer
  (\cref{chap:hyperbolic_family}).
\end{lemma}

\begin{proof}
  \uses{lem:arc_trajectory_diff_bound, lem:gronwall_L1_drive}
  \leanok
  \cref{lem:arc_trajectory_diff_bound} feeds the Grönwall drive lemma
  \cref{lem:gronwall_L1_drive}.
\end{proof}

\begin{lemma}\leanok
[a confined solution realizes $\kappa$ at $\varepsilon=-1$]
  \label{lem:arc_solution_realizes}
  \lean{Gluck.Hyperbolic.arcSolution_realizes}
  \uses{def:arc_angle_speed, def:spaceform_realizes}
  If $(z,\varphi)$ solves the \emph{true} (unclamped) system and stays confined
  ($\|z(\sigma)\|<1$ for all $\sigma$), then
  $\mathrm{Realizes}\,(-1)\,z\,\kappa$ with tangent angle $\varphi$. No new
  predicate is needed: the space-form realization predicate of
  \cref{chap:spaceform_core} is already parametrization-flexible ($\varphi$ may
  be non-monotone, the bracket may be negative).
\end{lemma}

\begin{proof}
\leanok
  $z$ is $C^1$ with $\|z'\|=1\ne0$, confined by hypothesis, and the gauge-speed
  relation is exactly the $\varphi$-ODE after clearing the positive metric
  factor.
\end{proof}

\begin{definition}\leanok
[constant-curvature arc model]
  \label{def:arc_model_const}
  \lean{Gluck.Hyperbolic.arcModelConst}
  \uses{def:arc_angle_speed}
  For $\kappa\equiv K$ the reconstruction through $(z_0,\varphi_0)$ is the
  explicit Euclidean circular arc
  $z(\sigma)=z_0-r\,ie^{i\varphi_0}(e^{i\sigma/r}-1)$,
  $\varphi(\sigma)=\varphi_0+\sigma/r$, of radius
  $r=(1-\|z_0\|^2)/\bigl(2(K+\langle z_0,ie^{i\varphi_0}\rangle)\bigr)$
  (\texttt{arcModelRadius}). While confined it solves the true
  system (\texttt{arcModelConst\_solves}), hence coincides with
  \texttt{arcFlow} via flow uniqueness; this closed form powers all
  family-stage computations.
\end{definition}

\begin{theorem}\leanok
[Poincaré--Miranda on a rectangle]
  \label{thm:poincare_miranda_rect}
  \lean{Gluck.Hyperbolic.poincareMiranda_rect}
  \uses{thm:existence_of_zero, def:winding_number_c}
  A continuous $G=(G_1,G_2):[a_1,a_2]\times[b_1,b_2]\to\mathbb R^2$ with each
  component sign-definite on its pair of faces ($G_1\le0$ left, $G_1\ge0$
  right; $G_2\le0$ bottom, $G_2\ge0$ top) has a zero in the rectangle. The 2-D
  intermediate value theorem; absent from mathlib, proved project-side. It is
  the topological engine of the closing crux: the quarter landing is a genuine
  2-D shooting problem in the mirror height $h$ and the free window length $L$
  — unlike the decoupled Euclidean $\varphi'=\kappa$, the coupled
  $\varphi$-turning is not automatic, and H\textsuperscript{2} has no rescaling
  to normalize $L$ away.
\end{theorem}

\begin{proof}
  \uses{thm:existence_of_zero}
  \leanok
  Via the project's planar degree principle. Rescale the rectangle to the unit
  disk, push $G$ forward to $F:\mathbb C\to\mathbb C$; the four faces make $F$
  nonvanishing on the boundary circle, and the boundary loop threads the four
  half-planes $\{\operatorname{Im}<0\}$, $\{\operatorname{Re}>0\}$,
  $\{\operatorname{Im}>0\}$, $\{\operatorname{Re}<0\}$ in cyclic order, so a
  four-arc computation in the winding API (\cref{def:winding_number_c}) gives
  winding number $1$; \cref{thm:existence_of_zero} supplies the interior zero.
  Strict faces first (\texttt{poincareMiranda\_rect\_strict}), then a
  compactness limit.
\end{proof}

\begin{definition}\leanok
[H\texorpdfstring{$^2$}{2} arc-length curvature function]
  \label{def:arclength_h2_curvature}
  \lean{Gluck.Hyperbolic.ArcLengthH2Curvature}
  \uses{def:arc_angle_speed, def:arclength_curvature}
  $\kappa:\mathbb R\to\mathbb R$ is an \emph{H\textsuperscript{2} arc-length
  curvature function} if there are a window length $L>0$ and a pair
  $(z,\varphi)$ globally solving the H\textsuperscript{2} arc-length system,
  confined ($\|z\|<1$), closing ($z(L)=z(0)$) with total turning $2\pi$
  ($\varphi(L)=\varphi(0)+2\pi$), $L$-periodic in $z$, and injective on
  $[0,L)$. Coupled H\textsuperscript{2} analogue of the Euclidean
  \cref{def:arclength_curvature} (Dahlberg's (1.1)--(1.3)); the essential
  difference is that $L$ is existentially quantified — co-constructed with the
  profile — because H\textsuperscript{2} admits no metric rescaling.
\end{definition}

\begin{lemma}\leanok
[floor-glued periodic extension of a window solution]
  \label{lem:window_solution_exposed}
  \lean{Gluck.Hyperbolic.windowSolution_exposed}
  \uses{def:arc_flow, def:arclength_h2_curvature,
        lem:injon_dahlberg_curve}
  Let $\kappa$ be continuous and $L$-periodic, and suppose the window flow
  $\Phi=\Psi(W_0,\cdot)$ on $[0,L]$ \emph{closes}
  ($\Phi(L)=(W_{0,z},W_{0,\varphi}+2\pi)$), is \emph{confined}
  ($\|\Phi_z\|\le R<1$) and satisfies the \emph{chord condition}
  $\int_t^\tau e^{i\varphi}\ne0$ on every proper sub-arc. Then the floor-glued
  periodic extension
  $Z(\sigma)=\Phi(\sigma-L\lfloor\sigma/L\rfloor)
  +\lfloor\sigma/L\rfloor\,(0,2\pi)$
  is a global solution of the true system, $L$-periodic in $z$, confined,
  closing with total turning $2\pi$, injective on $[0,L)$ — all the data of
  \cref{def:arclength_h2_curvature} at window $L$.
\end{lemma}

\begin{proof}
  \uses{def:arc_flow, lem:injon_dahlberg_curve}
  \leanok
  The gluing lemma (\texttt{periodic\_glue}) differentiates the extension
  across the seams: the endpoint match and the field-endpoint match
  ($\kappa(L)=\kappa(0)$, $e^{i(\varphi+2\pi)}=e^{i\varphi}$) align the
  one-sided derivatives, and $L$-periodicity of $\kappa$ transports the field
  along the floor shift. Confinement deactivates the clamp, so $Z$ solves the
  \emph{true} system. Injectivity is the chord argument
  (\texttt{injOn\_arcCurve}): the FTC bridge $z(b)-z(a)=\int_a^b e^{i\varphi}$
  turns equal values into a vanishing chord — a direct reuse of the
  Euclidean-in-disk argument (\cref{lem:injon_dahlberg_curve}), embeddedness
  being independent of the H\textsuperscript{2} metric.
\end{proof}

\begin{lemma}\leanok
[the converse at an explicit window, with the shift law]
  \label{lem:arclength_h2_converse_at}
  \lean{Gluck.Hyperbolic.arcLengthH2Converse_at}
  \uses{lem:arc_solution_realizes, def:is_simple_closed,
        def:spaceform_realizes}
  Given a confined, $L$-periodic-in-$z$, simple global solution $(z,\varphi)$
  of the H\textsuperscript{2} arc-length system for continuous $\kappa$ and
  window $L>0$, the linear window reparametrization $\chi(t)=(L/2\pi)t$
  produces $Z=z\circ\chi$ with: $\chi$ is $C^1$ with $\chi'>0$, the
  \emph{window shift law} $\chi(t+2\pi)=\chi(t)+L$, $Z$ is simple closed, and
  $\mathrm{Realizes}\,(-1)\,Z\,(\kappa\circ\chi)$. The shift law is the
  window-conjugation datum consumed by the degree-one analysis of
  \cref{chap:hyperbolic_converse} (the exact-profile upgrade).
\end{lemma}

\begin{proof}
  \uses{lem:arc_solution_realizes}
  \leanok
  \cref{lem:arc_solution_realizes} gives $\mathrm{Realizes}\,(-1)\,z\,\kappa$;
  the no-rescaling transport \texttt{spaceFormRealizes\_comp}
  (\cref{chap:spaceform_converse} infrastructure) carries it along $\chi$ to
  $\kappa\circ\chi$ — \emph{not} $\kappa$: there is no scaling step, unlike the
  Euclidean \cref{thm:arclength_converse}. The shift law is
  $\chi(t+2\pi)=(L/2\pi)t+L$. Simplicity: $\chi$ maps $[0,2\pi)$ bijectively
  onto $[0,L)$, transporting $L$-periodicity of $z$ to $2\pi$-periodicity of
  $Z$ and injectivity on $[0,L)$ to injectivity on $[0,2\pi)$.
\end{proof}

\begin{theorem}\leanok
[the H\texorpdfstring{$^2$}{2} arc-length converse, up to reparametrization]
  \label{thm:arclength_h2_converse}
  \lean{Gluck.Hyperbolic.arcLengthH2Converse}
  \uses{def:arclength_h2_curvature, lem:arclength_h2_converse_at,
        def:is_simple_closed, thm:arclength_converse}
  If $\kappa$ is continuous and an H\textsuperscript{2} arc-length curvature
  function, there are a simple closed curve $z$ and an orientation-preserving
  $C^1$ reparametrization $\psi$ with
  \[
    \texttt{IsSimpleClosed}\,z \;\wedge\;
    \mathrm{Realizes}\,(-1)\,z\,(\kappa\circ\psi).
  \]
  Honest H\textsuperscript{2} analogue of \cref{thm:arclength_converse}, with
  the metric-rescaling step replaced by the window reparametrization: $L$ is
  co-constructed and generically $\ne2\pi$, so an on-the-nose $2\pi$-convention
  realization of $\kappa$ is not available at this stage (an earlier exact
  conclusion was unsound); the reparametrization is removed later, for the
  concrete construction, by the degree-one analysis of
  \cref{chap:hyperbolic_converse}. No $2\pi$-periodicity of $\kappa$ is
  assumed — for a co-constructed arc-length profile it can fail.
\end{theorem}

\begin{proof}
  \uses{lem:arclength_h2_converse_at}
  \leanok
  Destructure the window existential and apply
  \cref{lem:arclength_h2_converse_at}, discarding the shift law.
\end{proof}

\begin{lemma}\leanok
[realization up to reparametrization]
  \label{lem:realizes_h2_of_reparam}
  \lean{Gluck.Hyperbolic.realizesH2_of_reparam}
  \uses{thm:arclength_h2_converse, def:arclength_h2_curvature,
        def:is_simple_closed}
  Let $\kappa$ be continuous and $2\pi$-periodic, and $\psi$ a $C^1$
  orientation-preserving circle map ($\psi(t+2\pi)=\psi(t)+2\pi$) such that
  $\kappa\circ\psi$ is an H\textsuperscript{2} arc-length curvature function.
  Then $\kappa$ is realized \emph{up to a further orientation-preserving $C^1$
  reparametrization} $\Psi$ by a simple closed H\textsuperscript{2} curve:
  $\texttt{IsSimpleClosed}\,z\wedge\mathrm{Realizes}\,(-1)\,z\,(\kappa\circ\Psi)$.
  Honestly up-to-reparam: the base converse closes at the co-constructed window
  $L'$, and an exact $2\pi$-realization would need $\psi$ to conjugate the
  $L'$-shift to the $2\pi$-shift, which the available $2\pi$-shift law cannot
  supply for $L'\ne2\pi$. The degree-one analysis of
  \cref{chap:hyperbolic_converse} removes the reparametrization for the
  concrete construction.
\end{lemma}

\begin{proof}
  \uses{thm:arclength_h2_converse}
  \leanok
  \cref{thm:arclength_h2_converse} applied to the continuous $\kappa\circ\psi$
  yields a simple closed $Z$ realizing $(\kappa\circ\psi)\circ\chi$ for its
  internal window reparametrization $\chi$; take $\Psi=\psi\circ\chi$, $C^1$
  with $\Psi'=(\psi'\circ\chi)\cdot\chi'>0$, and
  $(\kappa\circ\psi)\circ\chi=\kappa\circ\Psi$ definitionally.
\end{proof}
